The low-frequency wave multiplier is smooth off the origin. The following multiplier bound controls it by derivative estimates there.

\begin{Proposition}[Mihlin--H\"ormander multiplier bound]\label{ext:mihlin-multiplier}
Let $n\ge1$, $N=\lfloor n/2\rfloor+1$ and $1<p<\infty$. For a complex symbol $m\in C^N(\mathbb R^n\setminus\{0\})$ assume
\[
 A=\max_{|\alpha|\le N}\ \sup_{\xi\ne0}|\xi|^{|\alpha|}
             |\partial_\xi^\alpha m(\xi)|<\infty.
\]
For $f\in\mathcal S(\mathbb R^n)$ define $m(D)f=(m\widehat f)^\vee$ using
\[
 \widehat f(\xi)=\int_{\mathbb R^n}e^{-ix\cdot\xi}f(x)\,dx,\qquad
 g^\vee(x)=(2\pi)^{-n}\int_{\mathbb R^n}e^{ix\cdot\xi}g(\xi)\,d\xi .
\]
The value of $m$ at the origin is immaterial. Then
\[
 \|m(D)f\|_{L^p(\mathbb R^n)}
 \le C_{n,p} A\|f\|_{L^p(\mathbb R^n)}
 \qquad (f\in\mathcal S(\mathbb R^n)).
\]
By density $m(D)$ extends uniquely to a bounded operator on $L^p(\mathbb R^n)$ satisfying the same inequality for every $f\in L^p(\mathbb R^n)$.
Indeed the displayed derivative bounds imply the annular $L^2$ bounds
$\|\partial^\alpha m\|_{L^2(\{R<|\xi|<2R\})}\le C_n A R^{n/2-|\alpha|}$
used in \cite{r-32}, Theorem~6.2.7 and (6.2.14), p.~446.
\end{Proposition}
